% !TEX root = ../main.tex
\section{Accounting vs.\ Economic Profit}\label{thr:accounting-vs-economic-profit}

\paragraph{Definisjon.} \emph{Regnskabsmæssig profit} (accounting profit) er inntekter minus eksplisitte bokførte kostnader. \emph{Økonomisk profit} (economic profit) er inntekter minus \emph{alle} kostnader, inkludert \emph{opportunity costs} -- verdien av den beste alternative bruken av ressursene. En aktivitet med positiv regnskapsprofitt kan ha negativ økonomisk profitt hvis alternativkostnaden overstiger regnskapsmarginen. [SLIDE]

\subsection*{Forklaring}

BSZ kap.\ 8 bruker distinksjonen til å forklare to ting: [SLIDE]

\textbf{1. Korrekt beslutningsgrunnlag.} Ledere som kun ser på regnskapstall kan beholde aktiviteter som \emph{ødelegger verdi} fordi de ignorerer alternativkostnaden. Økonomisk profitt er det riktige beslutningskriteriet -- det fanger den reelle verdiskapningen utover hva ressursene ville tjent i sin beste alternative anvendelse. [BOK]

\textbf{2. Boundaries of the Firm.} Hvis en vertikal integrert enhet viser positiv regnskapsprofitt men negativ økonomisk profitt, bør aktiviteten \emph{outsources}: markedet kan utføre den billigere. Eksemplet er et oljeselskap som raffinerer selv -- raffineringsmarginene dekker regnskapskostnadene, men opportunity-kosten ved å binde kapital i et raffineri (i stedet for å reinvestere i leting) gjør at økonomisk profitt er negativ. [SLIDE]

\textbf{Formell sondring:}
\[
\pi_{\text{accounting}} = R - C_{\text{explicit}}, \qquad
\pi_{\text{economic}} = R - C_{\text{explicit}} - C_{\text{opportunity}}
\]
Forskjellen, $C_{\text{opportunity}}$, inkluderer: (i) alternativ avkastning på investert kapital, (ii) verdi av eierens tid/kompetanse i annen anvendelse, (iii) alternativ salgspris for innsatsfaktorer i markedet. [BOK]

\subsection*{Kjerneforutsetninger}
\begin{itemize}
\item Et velfungerende marked eksisterer for innsatsfaktorene, slik at opportunity cost kan prises.
\item Ledere har (eller bør ha) informasjon om alternativavkastning.
\item Regnskapssystemet dekker eksplisitte kostnader korrekt (avskrivninger, renter, lønninger).
\end{itemize}

\subsection*{Muligheter}
\begin{itemize}
\item Avslører om vertikal integrasjon faktisk skaper verdi eller bare maskerer dårlige alternativer.
\item Gir normativt kriterium for \emph{make-or-buy}: outsource hvis $\pi_{\text{economic}}<0$.
\item Kobler til EVA/residual income-modeller i \thref{responsibility-centers}{investment centers}.
\end{itemize}

\subsection*{Begrensninger og kritikk}
\begin{itemize}
\item Opportunity cost er vanskelig å estimere presist -- avhenger av kontrafaktiske scenarioer.
\item I fravær av likvide markeder (f.eks.\ for spesifikke aktiva) er alternativverdien uklar.
\item Regnskapsprofitten er det eneste som rapporteres eksternt; økonomisk profitt krever intern analyse.
\item Strategiske hensyn (kontroll, sikkerhet, market power) kan rettferdiggjøre aktiviteter med negativ $\pi_{\text{economic}}$ -- modellen er nødvendig men ikke tilstrekkelig.
\end{itemize}

\subsection*{Modell / illustrasjon}

\begin{figure}[H]
  \centering
  \includegraphics[width=0.7\linewidth]{BSZ_kap_8_2026/by_slide/slide_011/BSZ_kap_8_2026_slide_011_asset_01_image6.png}
  \caption{BP Retail Gasoline Stations: regnskapsmessig vs.\ økonomisk profitt. Ved \$3{,}60/gallon er regnskapsmessig profitt positiv (\$1{,}2M) men økonomisk profitt negativ (-\$400k) fordi opportunity cost for bensin solgt i eget nettverk overstiger marginen. Eksemplet viser at BP bør outsource/selge detaljistoperasjonen. [SLIDE -- BSZ kap.\ 8, Table 8.1]}
\end{figure}

\subsection*{Eksempler}
\begin{itemize}\setlength{\itemsep}{2pt}
\item \textbf{BP gasoline stations (BSZ Table 8.1):} Regnskapsprofitt \$1{,}2M vs.\ økonomisk profitt -\$400k ved lavere bensinpris -- illustrerer at detaljistleddet ødelegger verdi. [SLIDE]
\item \textbf{Norsk kommune og egen IT-drift:} Kommunen viser lave IT-kostnader i regnskapet, men inkluderer ikke alternativkostnaden av IT-personalets tid og kommunens forhandlingsmakt i et konkurranseutsatt marked. Outsourcing kan gi positiv økonomisk profitt selv om regnskapskostnadene stiger. [EKS]
\end{itemize}

\subsection*{Kobling til søyle og andre teorier}
Søyle: \STwo\ (korrekt prestasjonsmåling krever riktig profittbegrep). Relaterte: \thref{boundaries-of-the-firm}{Boundaries of the Firm}, \thref{responsibility-centers}{Responsibility centers} (EVA i investment centers), \thref{transfer-pricing}{Transfer pricing} (intern prising reflekterer opportunity cost), \thref{vertical-integration-outsourcing}{Vertikal integrasjon og outsourcing}.

\paragraph{Brukt i spørsmål:} Q11.
